% !TeX program = xelatex
% UTF-8 transcription; original wording, figures and abbreviations retained.
% pic01: 20261016_044537060.jpg; pic02: 20261016_044542025.jpg;
% pic03: 20261016_044544962.jpg; pic04: 20261016_044549691.jpg;
% pic05: 20261016_044552691.jpg.
% Bracketed red text identifies uncertain handwriting or text cut off in the photo.
% Standalone XeLaTeX document, also usable via input in a Xepersian document.
\ifdefined\settextfont
  \def\AtopicDermatitisStandalone{0}
\else
  \def\AtopicDermatitisStandalone{1}
  \documentclass[12pt,a4paper]{article}
  \usepackage[margin=2cm]{geometry}
  \usepackage{amsmath,amssymb,xcolor,enumitem}
  \usepackage{xepersian}
  \IfFontExistsTF{Amiri}{\settextfont{Amiri}}{\settextfont{FreeSerif}}
  \setlatintextfont{FreeSerif}
\fi
\providecommand{\en}[1]{\lr{#1}}
\newcommand{\adunclear}[1]{\textcolor{red!65!black}{[خوانش نامطمئن: #1]}}
\newcommand{\adcut}{\textcolor{red!65!black}{[ادامه خارج از کادر عکس]}}
\newcommand{\adpage}[1]{\par\bigskip\noindent\textbf{صفحهٔ #1 ـ اگزما آتوپیک ـ ۴۱}\par\smallskip\noindent موضوع: اگزما آتوپیک ـ ۴۱ \hfill تاریخ: / /\par\medskip}
\ifnum\AtopicDermatitisStandalone=1 \begin{document}\fi

\adpage{۱}
\noindent\en{Rook's Dermatology (10th)}

یک وضعیت خارش‌دار، عودکننده، به‌طور مزمن و التهابی پوست است که عموماً در اوایل کودکی آغاز می‌شود (معمولاً قبل از ۲ سالگی).

\begin{itemize}
  \item ویژگی بارز: التهاب (قرمزی در پوست‌های روشن و افزایش التهاب در پوست‌های \adcut{} تیره).
  \item تغییرات سطح پوست: \en{Fine scaling}، اکسکوریشن، لیکنیفیکاسیون.
  \item معمولاً در زمینهٔ خارش مداوم؛ ادم، پاپولیشن، اوزینگ (\adunclear{تراوش؛ انتهای سطر در لبهٔ عکس}).
  \item عموماً سطوح فلکسورال را درگیر می‌کند.
  \begin{itemize}
    \item درگیری غیر فلکسورال در شیرخواران شایع است و سیاه‌پوستان \adcut
  \end{itemize}
\end{itemize}

\en{AD} معمولاً با سایر وضعیت‌های آتوپیک همراهی دارد.

\subsection*{تظاهرات بالینی}
خارش، اریتم ماکولار، پاپول/پاپولووزیکول، ضایعات اگزماکوز ـ به‌ویژه \adunclear{کراتوتیک}، لیکنیفیکاسیون، هایپر و هایپوپیگمانتاسیون \adcut، خشکی پوست، عفونت ثانویه.

\subsection*{کرایتریا تشخیص}
کودک باید حتماً پوست خارش‌دار طی ۱۲ ماه گذشته \adcut؛ حداقل ۳ مورد:
\begin{enumerate}
  \item سن شروع زیر ۲ سال.
  \item سابقهٔ درگیری فلکسورال.
  \item سابقهٔ پوست خشک.
  \item سابقهٔ سایر بیماری‌های آتوپیک.
  \item درماتیت فلکسورال قابل مشاهده.
\end{enumerate}
\noindent حاشیهٔ معیار پنجم: آسم در سن زیر ۴ سال \adcut؛ اقوام درجهٔ ۱ هم حساب \adcut

\adpage{۲}
\begin{itemize}
  \item بیماران \en{AD} می‌توانند سطح نرمال \en{total IgE} داشته باشند.
  \item \en{AD} \adunclear{لزوماً} مرتبط با یک آلرژن خاص است.
\end{itemize}

\begin{latin}
\noindent AD $\longrightarrow$ Allergen-specific IgE $\uparrow$ (40--90\% of cases)\\
\hspace*{2em}(Extrinsic AD)\par\medskip
\noindent AD $\longrightarrow$ N/L IgE level (50\% $<$ of cases)\\
\hspace*{2em}(Intrinsic AD or non-atopic AD)
\end{latin}
% The comparison sign above follows the handwritten English line literally.

\begin{itemize}
  \item در طول زمان و افزایش عمر موارد \en{E-AD} شایع‌تر می‌شوند (تا ۹۵٪).
  \item شیوع \en{AD} در ایران ۲٪ و در ژاپن ۱۶٪ است؛ \en{$M<F$}.
  \item شیوع سالانهٔ آن حدود ۲۰٪ است.
  \item شیوع آن در میان‌سالی و پیری نیز افزایش می‌یابد.
  \item در \en{AD}: عوامل ژنتیکی و عوامل محیطی $\leftarrow$ تغییر عملکرد پوست $\leftarrow$ التهاب پوست $\leftarrow$ \en{AD}.
  \item موارد شدید \en{AD} نامعمول هستند و تنها ۵٪ موارد را تشکیل می‌دهند.
  \item ۲–۱ وعده میوهٔ تازه در هفته احتمال \en{AD} را کم می‌کند.
  \item ۳ بار فست‌فود در هفته احتمال \en{AD} را زیاد می‌کند.
\end{itemize}

\adpage{۳}
\begin{itemize}
  \item مصرف ماهی طی بارداری ریسک \en{AD} طی ۵ سال اول زندگی را \adunclear{کم می‌کند؛ پایان سطر بریده است}.
  \item مکمل‌های تغذیه‌ای نقش خاصی ندارند.
  \item تغذیه با شیر مادر احتمال \en{AD} را کاهش می‌دهد.
  \item چاقی هم ریسک بروز و هم شدت \en{AD} را افزایش می‌دهد.
  \item مصرف دخانیات در پدر و مادر ریسک \en{AD} را تغییر نمی‌دهد.
  \item آلودگی هوا به \en{NO\textsubscript{2}} تغییر نمی‌دهد.
  \item مهدکودک در ۲ سال اول زندگی ریسک را بالا می‌برد.
  \item حیوان خانگی در سال‌های اولیهٔ زندگی ریسک را کم می‌کند (خصوصاً سگ).
\end{itemize}

\subsection*{مارش آتوپیک}
\begin{itemize}
  \item پیشروی \en{AD} به سایر بیماری‌های آتوپیک:
  \begin{itemize}
    \item رینیت آلرژیک، آلرژی غذایی، آسم.
    \item تنها در ۷٪ موارد \en{AD} رخ می‌دهد.
  \end{itemize}
  \item \en{IgE} نقش جزئی در \en{AD} داشته چرا که مهار آن کارساز نیست $\leftarrow$ اومالیزومب.
  \item در \en{AD} سطح ائوزینوفیل و نوتروفیل بالاست.
  \item ۲۱٪ بیماران \en{AD} آلرژی غذایی دارند (۴٪ در نرمال جمعیت).
  \item ارتباطی بین درماتیت آلرژیک تماسی (\en{ACD}) و \en{AD} وجود ندارد.
\end{itemize}

\adpage{۴}
\begin{itemize}
  \item خارش در \en{AD}: سیگنالینگ غیر هیستامینرژیک دارد.
  \begin{itemize}
    \item با جزء التهابی \en{AD} مرتبط است: \en{IL-31}، \en{TSLP}.
    \item به همین جهت با داروهای ضدالتهابی رفع می‌شود.
    \item سیکلوسپورین و مهارکننده‌های \en{JAK}.
    \item سطح \en{IL-31} با شدت خارش مرتبط است.
    \item سطح \en{IL-31} به دنبال سوپرآنتی‌ژن استافیلوکوک به‌شدت بالا می‌رود.
  \end{itemize}
\end{itemize}

\subsection*{تظاهرات بالینی \en{AD}}
\begin{itemize}
  \item سن شایع زیر ۲ سال است؛ خصوصاً ۶–۲ ماهگی.
  \item علامت اصلی خارش، به همراه راش مزمن \en{fluctuating} است.
  \item اختلال خواب به علت خارش؛ اضطراب و تخریب پوستی.
  \item محرک خارش: گرما، عرق، حمام، ورزش، استرس، لباس پشمی.
  \item ویژگی ضایعهٔ اگزما حاد: \adunclear{سنتریپتال}، حاشیه نامنظم، ملتهب، اندکی برجسته؛ بیشتر سطوح فلکسورال.
  \item اگزما مزمن: پوست کلفت می‌شود، لیکنیفیکاسیون، هایپرپیگمانتاسیون.
  \item کراست‌های زرد $\leftarrow$ عفونت ثانویهٔ باکتریال.
  \item پوست خشک است به دلیل تبخیر زیاد، به دلیل اختلال ساختاری پوست.
  \item در شیرخواران: بیشتر در صورت / سطوح فلکسورال.
\end{itemize}

\adpage{۵}
\begin{itemize}
  \item در کودکان: بیشتر سطح فلکسور، آرنج و زانو / سطوح کناری گردن / مچ دست و مچ پا.
  \item پچ‌های \adunclear{دیسکرومیه} ممکن است دیده شود.
  \item درگیری دست که گاهی با وزیکولیشن همراه است ممکن است دیده شود (\en{Pompholyx}).
  \item اتواینوکولیشن \adunclear{دلیل اول} عفونت باکتری و ویروسی را \adcut
  \item حاشیهٔ سطر قبل: \adunclear{بیشتر ...؛ ادامه در لبهٔ چپ عکس بریده است}.
  \item در بزرگسالان:
  \begin{itemize}
    \item بیشتر سطوح فلکسورها و دست.
    \item لیکنیفیکاسیون شایع‌تر است.
  \end{itemize}
\end{itemize}

\subsection*{منج درماتیت آتوپیک}
\begin{itemize}
  \item باید براساس شدت بیماری و نیاز بیمار باشد.
  \item حتماً شدت بیماری قبل شروع درمان ارزیابی و مکتوب گردد.
  \item درمان خط اول:
  \begin{itemize}
    \item کاهش مواجهه با تحریک‌کننده‌ها.
    \item حمام با مواد مرطوب‌کننده.
    \item استفادهٔ مداوم از مرطوب‌کننده $\leftarrow$ حتی در رمیژن / در همهٔ بیماران \adcut
    \item سرکوب التهاب با کمک \en{TCS} $\leftarrow$ \adunclear{در مرحلهٔ شدید در بیماران هم اتفاق} \adcut
  \end{itemize}
  \item معمولاً نیازی به تغییر رژیم غذایی نیست.
  \item آلرژن‌های تنفسی هم ممکن است دخیل باشند.
\end{itemize}

% Printed notebook labels on each source page: موضوع ; تاریخ / / .
% Handwritten page numbers 1--5 are preserved in the page headings.
\ifnum\AtopicDermatitisStandalone=1
  \def\AtopicDermatitisFinish{\end{document}}
\else
  \let\AtopicDermatitisFinish\relax
\fi
\AtopicDermatitisFinish
